\section{From Local Ordering to Connected Supervision}
\label{sec:theory}

\eqref{eq:lvalue} sums \emph{within-group} variances, so the question worth asking is what a model
that drove it to zero would have been shown; the answer is precise enough to use as a design rule.
It concerns $q_\theta$ rather than $p_\theta(x \mid c)$, and the zero set of the ideal residual
condition rather than the implemented update's trajectory. Write
$w(x) = \log q_\theta(x \mid c) + E(c,x)/kT$, let $\rho_j$ be the reference measure of group $j$ with
support $U_j$, and join two groups by an edge of the \emph{overlap graph} $\mathcal{G}_c$ when
$\rho_j$ and $\rho_{j'}$ are not mutually singular, a property of the measures for which
$U_j \cap U_{j'} \neq \emptyset$ is necessary but far from sufficient (\Cref{def:overlap}).

\begin{proposition}[What zero group-wise variance fixes]
\label{prop:ident}
For strictly positive group weights $\pi_j$ and finite second moments, let the population objective $\sum_j \pi_j \Var_{\rho_j}[w]$ vanish. Then \emph{(i)} $w \equiv b_j$
$\rho_j$-almost surely on $U_j$ for a scalar $b_j$: the residual is fixed to one constant per group,
and no normaliser is needed, since $Z(c)$ enters $w$ only through that constant.
\emph{(ii)} $b_j = b_{j'}$ along every edge of $\mathcal{G}_c$, so the constants agree within each
connected component; connectivity is sufficient; in the unrestricted measurable residual class a disconnected graph permits different constants. Continuity or model-class restrictions can impose further relations. \emph{(iii)} If the $\rho_j$
have positive Lebesgue densities and $\mathcal{G}_c$ has $C \ge 2$ components $U^{(1)},\dots,U^{(C)}$
with finite positive restricted partition integrals and carrying all the model's mass, the zero-loss densities form a $(C-1)$-parameter family
$q_\theta = a_i e^{-E/kT}$ on $U^{(i)}$ with $\sum_i a_i Z^{(i)} = 1$, in which every vector of
component occupancies in the open simplex is attained: the \emph{relative} mass of separated
components is not identified.
\end{proposition}

\Cref{thm:l2g} states this in full, with the connected case and its two further hypotheses, and
\Cref{app:l2g} proves it. Our runs sit in the empirical reading (\Cref{cor:empirical}): $\rho_j$ is
the empirical measure on a group's $K$ sampled geometries, zero loss is $K-1$ scalar equations at
those points and nothing about unsampled ones, and parts (ii) and (iii) survive only as an
identifiability statement.

\paragraph{What our design supervises.} Each parent contributes exactly one group, and distinct
parents draw displacements independently from continuous distributions, so with probability one no
two groups share a sampled point: $\mathcal{G}_c$ is a set of isolated vertices and only part~(i)
applies (\Cref{cor:basins}). Part~(i) describes an ideal zero residual, which the experiments do not achieve --- the archived readout is not a validated $\log q_\theta$ --- and its complement is equally exact, the cloud's offset being free and the
relative probability of distinct basins untouched. Disconnection rests on mutual singularity rather
than on one group per composition, which our pool does not satisfy: $31\%$ of its parents sit on a
repeated composition, and those clouds are still disconnected (\Cref{app:l2g,app:data}). Nothing
above predicts the \emph{magnitude} of a correlation or the calibration of $\log q_\theta$
(\Cref{app:scope}), which are separate quantities measured separately in \Cref{sec:scale}.

\paragraph{The design criterion.} The analysis therefore turns an ambiguity into a design choice.
Connectivity is a property of the supervision graph, not of the model or the potential: local groups
suffice for local ordering, which is what we build and measure, while controlling relative basin
mass requires either a connected supervision graph --- groups tiled along a path between basins, so
part~(ii) propagates one constant across them --- or an additional global objective. This paper
measures local ordering; \Cref{app:nonequilibrium} specifies a global path-weight correction
as a separate development branch, without claiming established molecular benefit.
